\documentclass[11pt]{article}

\usepackage{amsmath,amssymb,amsthm}
\usepackage{booktabs}
\usepackage[margin=1in]{geometry}
\usepackage{hyperref}
\usepackage{url}
\setlength{\emergencystretch}{3em}

\newtheorem{proposition}{Proposition}
\newtheorem{remark}{Remark}

\title{Keep the Angle, in Practice:\\
A Substrate-Agnostic Manifold Diagnostic and Its Failure Modes}
\author{Andrew H. Bond\\
\small Department of Computer Engineering, San Jos\'e State University\\
\small San Jos\'e, CA 95192 USA \quad \texttt{andrew.bond@sjsu.edu} \quad ORCID 0009-0003-2599-6158}
\date{}

\begin{document}
\maketitle

\begin{abstract}
The companion theory paper (\emph{Keep the Angle}, Paper~I) proves that the angular
coordinate of a low-Laplacian spectral embedding preserves geodesic geometry while the
radial coordinate is degree/density nuisance. This paper is the empirical study of that
basis as a \emph{practical, substrate-agnostic diagnostic}. On graph families spanning
random-geometric tori and king-lattices, manifold-embedding samples (swiss-roll),
Lorentzian causal sets, and hypergraph-rewriting (Wolfram-model) substrates, the angular
map recovers graph geodesics wherever the substrate carries genuine low-dimensional
Riemannian geometry, independent of how the graph was generated: on a $2$-torus the
angular Spearman correlation is $\approx0.92$ and flat in the mode count, while the full
commute embedding decays; in a bake-off it beats every closed-form spectral competitor
($0.894\pm0.008$ vs.\ $0.60$--$0.77$ on the $2$-torus over five independent draws). We
show the same map is a \emph{diagnostic}: a two-factor small-world/dimension screen
separates genuine manifolds from destroyed or non-manifold geometry, angular fidelity
declines monotonically with measured intrinsic dimension across five Wolfram-model rule
families (an exploratory trend, $\text{angle-}\rho\approx1.09-0.157\,d$), and under
continuum refinement the basis \emph{sharpens} on sampled manifolds while \emph{decaying}
on emergent hypergraph geometry---returning opposite verdicts for sampled and emergent
substrates. Two pre-registered physical predictions are honest negatives. Everything here
is pre-registered and reproducible; the theory it rests on is Paper~I.
\end{abstract}

\section{Introduction}
\emph{Keep the Angle} (Paper~I) establishes, as theorems, that row-normalizing a
low-mode spectral embedding---projecting each node to the unit sphere and keeping only
its direction---recovers the graph's geodesic geometry, because the discarded radial
coordinate is degree/density nuisance (von~Luxburg-degenerate for the full unnormalized
commute radius at $d\ge3$) while the angular coordinate carries the metric. That paper is
about \emph{why} the angle works: an exact normalization identity, a deterministic
angular-transfer theorem, a conditional graph-to-continuum result, the unconditional flat
torus and sphere instances, a uniform lower angular bound on flat tori, and the filter
dichotomy (heat uniform, commute one-sided, plain rank-collapsing).

This paper is about \emph{whether, where, and how well} the angle works in practice, and
what it does when the geometry is not there. Our thesis is that the angular basis is a
\emph{substrate-agnostic manifold diagnostic}: on any graph that carries genuine
low-dimensional Riemannian geometry it recovers that geometry, independent of graph
origin; on graphs that do not, its failure is \emph{predictable and informative} rather
than a silent artifact. We support this with a controlled measurement program.

\paragraph{Contributions.}
(1) A geometry-recovery benchmark across five substrate families with matched controls
(\S\ref{sec:recovery}), and a bake-off in which the angle beats every purely spectral
distance built from the same spectrum (\S\ref{sec:bakeoff}).
(2) A study of \emph{emergent} geometry from hypergraph rewriting---Wolfram-model
substrates that are geometry with no latent coordinates---including an exploratory
fidelity-versus-dimension trend and temporal stability across substrate evolution
(\S\ref{sec:emergent}).
(3) The angular basis as an explicit \emph{diagnostic}: a two-factor screen that
separates manifolds from non-manifolds, and a continuum-refinement test that dissociates
sampled from emergent geometry (\S\ref{sec:diagnostic}).
(4) Two honest negatives (\S\ref{sec:neg}) and a scoped physical interpretation
(\S\ref{sec:physical}).
We do not re-prove Paper~I; theorem references below point there.

\section{The angular basis and measurement protocol}\label{sec:protocol}
Given a graph with normalized Laplacian $L=I-D^{-1/2}AD^{-1/2}$, eigenpairs
$(\lambda_k,v_k)$, the low-mode embedding places node $i$ at
$\Psi_i=(v_1(i)/\sqrt{\lambda_1},\dots,v_m(i)/\sqrt{\lambda_m})$, with radius
$r_i=\lVert\Psi_i\rVert$ and angle $\hat u_i=\Psi_i/r_i$. We score the angular geometry
by the Spearman rank correlation (``angle-$\rho$'') between the chordal angular distances
$\lVert\hat u_i-\hat u_j\rVert$ and true graph geodesics over a fixed anchor set, against
three controls: \emph{random} ($m$ random modes), \emph{magnitude} (radius only), and
\emph{Isomap} (classical MDS of the geodesics, an oracle given the geodesic matrix).
Intrinsic dimension $d$ is measured by three estimators (ball-growth, spectral dimension,
effective rank); low inter-estimator spread certifies a manifold. Protocol details, draw
counts, and reproducibility tiers mirror Paper~I; all runs are pre-registered
(\texttt{PREREG\_RUNG3/4}) and released in the repository,
\url{https://github.com/ahb-sjsu/the-angular-observer}.

\section{Geometry recovery across substrates}\label{sec:recovery}
On a $2$-torus ($n=2000$) the polar decomposition separates cleanly: the angle carries
the geometry (flat across the gap-closed cutoffs $4,8,12,20$; e.g.\ the $m{=}20$ value
is $0.891\pm0.009$ over five independent draws), the magnitude is uninformative, and
the full commute embedding is
worse and decays with $m$. A random-mode basis reconstructs geodesics at $\rho\approx0$.

Across substrate families (Table~\ref{tab:substrates}), every graph with genuine low-$d$
Riemannian geometry shows the effect, independent of origin; the metric fails exactly
where the substrate fails the manifold criterion---the Lorentzian causal set (whose
dimension estimators scatter over $2.9$--$10.7$) and the rewired lattice---the predicted
failures, not exceptions. Two caveats keep the claim honest. On the two flat substrates
classical MDS beats the angle (it wins where the manifold embeds near-isometrically in
low-dimensional Euclidean space); the angle wins where topology obstructs such an
embedding (the torus) and, unlike MDS, is uniform in $m$ and $n$. And the swiss-roll is a
manifold-embedding case: the basis is identical whether the manifold comes from
hypergraph rewriting or from points sampled on a continuous manifold.

\begin{table}[h]\centering
\caption{Geometry recovery across substrate families ($m{=}20$, $150$ anchors).}
\label{tab:substrates}
\begin{tabular}{@{}lcccc@{}}
\toprule
substrate & meas.\ $d$ & angle-$\rho$ & random & Isomap$_{d=2}$ \\
\midrule
2-torus (reference)         & 2               & 0.89 & 0.01 & 0.67 \\
king-lattice                & 2               & 0.82 & 0.01 & 0.95 \\
swiss-roll (embedding)      & 2               & 0.93 & 0.01 & 0.99 \\
emergent $R_{3D}$ (Wolfram) & 2.07            & 0.82 & 0.01 & 0.81 \\
causal set (Lorentzian)     & $2.9$--$10.7$   & 0.57 & 0.01 & 0.55 \\
small-world (destroyed)     & $3.6$--$4.3$    & 0.45 & 0.00 & --- \\
\bottomrule
\end{tabular}
\end{table}

\subsection{A bake-off against named spectral distances}\label{sec:bakeoff}
Scoring every closed-form spectral distance on the identical graph and anchor set
(Table~\ref{tab:bakeoff}), the angle beats each purely spectral competitor, and the
margin widens in $d{=}3$ where commute degeneracy bites harder. The most informative row
is the amplified commute distance of von~Luxburg, Radl and Hein, designed precisely to
undo commute degeneracy by \emph{subtracting} the $1/k_i+1/k_j$ term: it barely improves
on raw commute and stays far below the angle, because subtracting the degenerate radial
term is not the same as \emph{dividing it out}, which is what row-normalization does
(Paper~I). Isomap wins only because it is handed the geodesics---an oracle ceiling, not a
competitor; the angle recovers $0.83$--$0.89$ of it from the spectrum alone. Angle scores
replicate across five independent graph draws: $0.894\pm0.008$ ($2$-torus, min--max
$0.886$--$0.907$) and $0.831\pm0.007$ ($3$-torus, $0.818$--$0.836$).

\begin{table}[h]\centering
\caption{Bake-off: Spearman $\rho$ vs.\ graph geodesics; angle scores are five-draw means.}
\label{tab:bakeoff}
\begin{tabular}{@{}lcc@{}}
\toprule
distance & $2$-torus & $3$-torus \\
\midrule
\textbf{angle (ours)}                       & \textbf{0.894} & \textbf{0.831} \\
full commute / resistance                   & 0.601 & 0.360 \\
biharmonic                                  & 0.735 & 0.515 \\
diffusion distance (best $t$)               & 0.771 & 0.631 \\
amplified commute                           & 0.628 & 0.364 \\
\midrule
\textit{Isomap, given the geodesics}        & \textit{0.971} & \textit{0.952} \\
\bottomrule
\end{tabular}
\end{table}

\subsection{Which weighting carries the angle, and density normalization}\label{sec:ablation}
The direction depends on the per-mode weight. Plain eigenmaps ($w_k{=}1$) give an angle
that \emph{collapses} with mode count (angle-$\rho$ falls from $0.554\pm0.013$ at
$m{=}10$ to $0.081\pm0.012$ at $m{=}20$ on the torus, five draws), because equal
weighting lets high modes swamp the direction; the commute weight
$1/\sqrt{\lambda_k}$ is flat instead ($0.918\pm0.006\to0.896\pm0.006$), as is
diffusion. Degree correction is \emph{invisible} to the angle: dividing or multiplying
each row by $\sqrt{k_i}$ leaves angle-$\rho$ unchanged to three digits, confirming the
degree lives in the radius, not the angle (Paper~I). Under strongly non-uniform sampling
the raw angle is \emph{not} sampling-invariant (angle-$\rho$ degrades from $0.93$ to
$0.79$ as a density contrast grows to $45\times$), because the graph Laplacian converges
to a density-weighted operator; the full Coifman--Lafon $\alpha{=}1$ normalization
restores it exactly (to $0.92$--$0.96$, flat across the contrast range), isolating density
into the radius and geometry into the angle.

\section{Emergent geometry from hypergraph rewriting}\label{sec:emergent}
Wolfram-model substrates are geometry with \emph{no latent coordinates}: a manifold, if it
appears, emerges from hypergraph rewriting rather than from sampling a continuous space.
Random arity-2 rules yield clean but only $\approx1.5$D manifolds; clean $2$D is rare but
not absent, reached both by a borrowed rule ($d=2.07\pm0.01$) and by from-scratch
evolutionary search. The angular basis holds on genuine $2$D emergent space
(Table~\ref{tab:substrates}).

\paragraph{An exploratory dimension--fidelity trend.} With only five independent rule
families, no ``law'' can responsibly be claimed. Across $20$ manifolds ($5$ rules $\times$
$4$ correlated seeds, $d\approx1$--$3.6$), angular fidelity falls with measured intrinsic
dimension, $\text{angle-}\rho\approx1.09-0.157\,d$, with a rule-clustered descriptive
interval $[-0.343,-0.094]$ four times wider than a naive point bootstrap; we read this as
a five-family exploratory monotonic trend, not an inferential interval. The mechanism is
the dimension-gated degeneracy of Paper~I; its rank-specific bound is consistent with, but
does not derive, the observed decay.

\paragraph{Stability across substrate evolution.} Robustness across substrate \emph{types}
is a statement about fixed graphs; the basis also persists as a substrate \emph{evolves}.
Embedding successive rewrite snapshots of $151$ qualifying rules on a fixed anchor set, the
angular structure tracks the evolving geometry at least as stably as the graph's own
geodesics on every rule resolvable at protocol scale ($143/143$ against a pre-registered
$80\%$ bar), and over the longest time lever angular persistence declines only gently with
emergent dimension ($-0.029$/dim), replicating independently in both arities.

\section{The angular basis as a diagnostic}\label{sec:diagnostic}
\paragraph{A two-factor screen.} A small-world coefficient $\omega$ and the
inter-estimator dimension spread separate genuine manifolds ($\omega\le-0.50$, spread
$\le0.51$) from strongly destroyed or non-manifold geometry ($\omega\ge+0.23$, spread
$\ge0.77$). The low angle-$\rho$ on the rewired lattice and the Lorentzian causal set is
thus small-world/non-manifold contamination, the predicted failure of the basis on
non-geometry, not a failure on genuine geometry. A king-lattice rewiring sweep bounds the
screen honestly: angle-$\rho$ collapses almost at once ($0.81\to0.47$ by $1\%$ rewiring)
while $\omega$ does not cross zero until $\approx5\%$, so angular fidelity is itself the
finer detector of geometric corruption in the transition band.

\paragraph{Continuum survival, and a dissociation.} Refining each substrate toward the
continuum with a sparse eigensolver ($N\in\{2,5,10,20,50\}\times10^3$) dissociates the two
substrate families. On random geometric graphs the angular fidelity is flat-to-rising in
$N$ (RGG $d{=}3$: $0.856\to0.912$, Spearman$(\rho,\log N)=+1.00$): refining the mesh makes
the low modes resolve the geometry \emph{better}. The emergent $R_{3D}$ manifold instead
\emph{decays} monotonically ($0.749\to0.654$, Spearman $-1.00$), falling below a
pre-registered floor. A sampled manifold refines to a clean continuum; an emergent
hypergraph, under this basis, does not. The angular basis is thus a diagnostic for whether
a substrate's fine structure stays manifold-like under refinement, returning different
verdicts for sampled ($2/3$ pass) and emergent ($1/1$ fail) geometry. One qualification:
the dissociation is stated at a fixed mode budget; a spectral-gap-matched rerun inverts the
direction, so ``clean geometry'' is a relation between mode-budget resolution and substrate
scale---the motivation for a dimension-adaptive basis.

\section{Honest negatives}\label{sec:neg}
Two pre-registered predictions failed. \textbf{Not secretly hyperbolic:} against a
binary-tree positive control, the tested substrates show no signature of latent hyperbolic
geometry. \textbf{No dynamical, energy-responsive curvature:} the strongest testable claim
of the observer-relative reading---that the coarse-grained curvature responds dynamically
to injected energy---is a null. We report both; neither undermines the basis-level and
diagnostic results above, and the physical interpretation is offered only as motivation.

\section{Physical interpretation and outlook}\label{sec:physical}
The decomposition ``separate scale from direction, keep the part that carries the task
geometry'' recurs beyond graph geometry: polar-coordinate KV-cache quantization splits key
vectors the same way, and the recurrence is evidence for the \emph{decomposition}, not for
always discarding the radius---the boundary being exactly whether the discarded radial
variable is nuisance for the metric preserved (the tangential-noncollapse condition of
Paper~I). For attention keys, whose geometry depends on per-channel scale, discarding the
radius destroys generation: scale is signal there, not nuisance. In the observer-relative
reading of emergent-geometry programs, the coarse-graining a bounded observer performs to
perceive geometry is a candidate for exactly this angular projection, but its strongest
dynamical claim is the null of \S\ref{sec:neg}, so we leave that door open. The operational
contribution stands independently: the angular basis is a substrate-agnostic diagnostic for
low-dimensional Riemannian structure, with predictable, informative failure modes. Open
directions: a dimension-adaptive basis that tunes mode-count and diffusion time against the
fidelity decay; substantially more independently selected rule families to test the
dimension trend; and evaluation against the latent manifold metric wherever coordinates
exist.

\paragraph{Reproducibility.} All pre-registration documents, experiment scripts, per-run
result files, rule sets, and random seeds are released in the accompanying repository,
\url{https://github.com/ahb-sjsu/the-angular-observer}. The theory these experiments test
is Paper~I (\emph{Keep the Angle}).

\begin{thebibliography}{9}
\bibitem{keeptheangle} A.~H.~Bond. \emph{Keep the Angle: A Geometry-Preserving Basis in
Spectral Embeddings} (Paper~I). Manuscript, 2026.
\end{thebibliography}

\end{document}
